%% =====================================================================
%% SHIELD — Self-Healing SDN framework paper (IEEEtran)   [S4.3]
%% Companion to main.tex (the FYDP closed-loop-controller paper). This
%% paper reframes the work around the SHIELD FRAMEWORK and its three
%% pillars (detect / resolve / learn), per FRAMEWORK_PLAN.md.
%%
%% COMPILES ON OVERLEAF (IEEEtran bundled) or any TeX Live:
%%   pdflatex shield -> bibtex shield -> pdflatex shield -> pdflatex shield
%% No TeX toolchain on the dev VM; build is S4.4 (Overleaf).
%%
%% HONESTY: every quantitative claim is tagged % [VERIFIED <result-doc>]
%% and traces to a reproducible artifact under docs/ + results/. Numbers
%% come from the HARDENED runs (p2 review M2/M3, p3 review M3/M4).
%% Shares paper/references.bib with main.tex.
%% =====================================================================
\documentclass[conference]{IEEEtran}
\usepackage{graphicx}
\usepackage{booktabs}
\usepackage{amsmath}
\usepackage{url}
\usepackage[hidelinks]{hyperref}

\begin{document}

\title{SHIELD: A Self-Healing SDN Framework that\\
Detects, Resolves, and Learns to Recover}

\author{\IEEEauthorblockN{Group B1}
\IEEEauthorblockA{Department of Computer Science and Engineering\\
United International University\\
Dhaka, Bangladesh}}

\maketitle

\begin{abstract}
Software-Defined Networking (SDN) centralizes control and makes recovery
programmable, yet most self-healing systems stop at a single stage---detecting a
fault, \emph{or} choosing a backup path, \emph{or} learning a policy---and evaluate
that stage in isolation. We present \textbf{SHIELD}, a closed-loop framework that
integrates all three into one control plane and evaluate each stage honestly on an
emulated Mininet/Open~vSwitch/Ryu testbed. (i)~\emph{Detection} uses a hybrid signal:
exact OpenFlow port-status for localization and normalized telemetry for
service-impact, localizing every injected hard fault ($19/19$, accuracy $1.000$)
across two topologies with zero congestion false alarms and a mean latency of
$0.21$~s. % [VERIFIED docs/p1_detection_results.txt]
(ii)~\emph{Resolution} scores surviving paths with a loss-aware objective
(latency, capacity, loss, hops); against loss-blind conventional recovery it wins in
$3/5$ live scenarios, ties a strong load-aware widest-path baseline on pure-bandwidth
faults, and \emph{beats} it when loss decorrelates from bandwidth (a gray failure on
a high-capacity link), while matching a ground-truth oracle under noisy telemetry.
% [VERIFIED docs/s2_5_results.txt]
(iii)~\emph{Learning} re-fits the resolver online on a non-stationary world where the
decisive signal (loss) is unobservable at decision time: over nine ``days'' and eight
seeds SHIELD-online significantly outperforms a frozen model ($p=0.0039$), tracks the
drift that collapses the frozen model, and \emph{matches} a principled adaptive
(sliding-window-greedy) baseline ($p=0.547$) without catastrophic forgetting.
% [VERIFIED docs/s3_4_results.txt]
Our contribution is the integrated, honestly-benchmarked framework---not a new
learning algorithm; indeed with accurate telemetry the heuristic resolver already
matches the oracle, and the learner's value is adaptation under partial observability
and drift. We release all components, topologies, fault-injection harness, and
evaluation scripts.
\end{abstract}

\begin{IEEEkeywords}
Software-Defined Networking, OpenFlow, self-healing, fault localization, path
selection, online learning, non-stationary, Ryu, Mininet, reproducibility.
\end{IEEEkeywords}

\section{Introduction}
Modern networks demand high availability, but a centralized SDN control plane
concentrates the risk of link and path failures. Static backup routes and manual
reconfiguration are often too slow or too rigid for dynamic conditions. A truly
self-healing controller must do three things in a loop: \emph{detect and localize}
a fault, \emph{resolve} it by choosing and installing a good recovery path, and
\emph{learn} from outcomes so that its choices improve as the network changes. Most
prior systems address only one of these and evaluate it in simulation.

We build and evaluate \textbf{SHIELD} (Self-Healing Integrated Empirical Learning
of Detection-resolution), a framework that closes this loop on real OpenFlow
actuation, and we hold each stage to an adversarial, leakage-free standard. Our
contributions are:
\begin{itemize}
  \item \textbf{An integrated self-healing framework} spanning detection,
  resolution, and learning over one Ryu/OVS control plane, with each stage exposed
  as a reusable component (Section~\ref{sec:arch}).
  \item \textbf{Hybrid fault localization} that combines exact OpenFlow port-status
  with telemetry, motivated by an empirical finding that pure-telemetry localization
  is infeasible; it localizes $19/19$ hard faults on two topologies with zero
  congestion false alarms (Section~\ref{sec:detect}).
  \item \textbf{A loss-aware resolver} benchmarked not only against loss-blind
  recovery (fixed-failover, shortest-path) but against a \emph{strong} load-aware
  widest-path baseline and a ground-truth oracle, isolating exactly where a
  multi-objective policy helps (Section~\ref{sec:resolve}).
  \item \textbf{An online learner} evaluated on a non-stationary world with an
  \emph{unobservable} decisive signal, against both a frozen ablation and a
  principled adaptive baseline, with a paired sign-flip permutation test over eight
  seeds (Section~\ref{sec:learn}).
  \item \textbf{An honesty-first methodology}: every number traces to a reproducible
  artifact, and we state plainly where a simpler method suffices---the framework,
  not the learning algorithm, is the contribution (Section~\ref{sec:threats}).
\end{itemize}

\section{Related Work}
\label{sec:related}
Learning-based control has been applied across SDN congestion management and network
fault recovery. For \emph{congestion}, Boussaoud \emph{et al.} combine a
classification agent with a deep-$Q$ decision agent to adapt routing and bandwidth in
a Mininet/Ryu testbed~\cite{boussaoud2025}. For \emph{link recovery}, Huang
\emph{et al.} use a DDPG agent to recompute load-balanced link weights after a
failure in a wide-area network~\cite{huang2025}, while Hammad \emph{et al.} train a
flow-based classifier that reduces intrusion-recovery time~\cite{hammad2023}. At the
infrastructure level, de~Souza~Neto \emph{et al.} keep control- and
management-plane communication alive under failure~\cite{neto2020}, and Li
\emph{et al.} formalize an RL life-cycle (``RLOps'') for Open~RAN~\cite{li2022rlops}.
A parallel line in power systems applies (D)RL to service
restoration~\cite{ferreira2019,jo2024}, and a recent review surveys AI-based
self-healing more broadly~\cite{kumar2023}.

Two gaps motivate SHIELD. First, most SDN studies stop at a \emph{single} decision
(a predicted congestion level, a recomputed weight vector) and evaluate it in
simulation, without integrating detection, path selection, and adaptation into one
loop that actually reprograms the data plane. SHIELD closes that loop end-to-end and
measures each stage on emulated OpenFlow. Second, several methods are labelled
reinforcement learning even though each decision is made from an independent
observation; we are explicit that our detector is contextual classification and that
our online stage is a non-stationary bandit, and we benchmark each against the
appropriate honest baseline rather than a strawman. Table~\ref{tab:related}
positions the closest work against the framework.

\begin{table*}[t]
\centering
\footnotesize
\caption{Positioning against representative learning-based SDN/self-healing work.}
\label{tab:related}
\begin{tabular}{@{}p{2.0cm}p{2.7cm}p{2.6cm}p{2.2cm}p{4.4cm}@{}}
\toprule
\textbf{Work} & \textbf{Problem} & \textbf{Method} & \textbf{Testbed} & \textbf{Scope vs. SHIELD} \\
\midrule
Boussaoud \emph{et al.} 2025~\cite{boussaoud2025} & SDN congestion control & MARL: classifier + DQN & Mininet/Ryu & Congestion only; no detect$\rightarrow$resolve$\rightarrow$learn loop or fault-recovery actuation. \\
Huang \emph{et al.} 2025~\cite{huang2025} & SDN faulty-link recovery & DDPG (link weights) & Simulation & Recomputes weights only; no emulated closed-loop, no online adaptation study. \\
Hammad \emph{et al.} 2023~\cite{hammad2023} & SDN intrusion recovery & ML flow classifier & SDN dataset & Security recovery, not link-fault localization + reroute. \\
de~Souza~Neto \emph{et al.} 2020~\cite{neto2020} & Control/mgmt-plane self-healing & Platform / architecture & Emulated & Keeps control-plane alive; not data-plane resolution/learning. \\
Ferreira 2019~\cite{ferreira2019}; Jo 2024~\cite{jo2024} & Power-grid restoration & (D)RL & Grid sim. & Different domain; informs restoration, not SDN. \\
\textbf{SHIELD (this paper)} & SDN link-fault self-healing & Hybrid detect $\rightarrow$ loss-aware resolve $\rightarrow$ online learn & Mininet/OVS/Ryu & Integrated loop; honest baselines at every stage (oracle, widest-path, adaptive bandit); reproducible. \\
\bottomrule
\end{tabular}
\end{table*}

\section{The SHIELD Framework}
\label{sec:arch}
SHIELD runs as a Ryu application over an Open~vSwitch data plane
(OpenFlow~1.3)~\cite{onf2012openflow,ryu} emulated in
Mininet~\cite{lantz2010mininet}. The control logic is split into a Ryu-free
``brain'' (so it is unit-testable off the data plane) and a thin Ryu actuator that
issues \texttt{OFPFlowMod} messages. Three cooperating stages form the loop:

\paragraph{Detection}
A \emph{telemetry collector} polls per-port statistics and normalizes each active
port to a vector $[\text{throughput},\text{loss},\text{delay}]$; a \emph{localizer}
consumes OpenFlow port-status events. Detection is deliberately \emph{hybrid}:
localization and hard-fault detection come from port-status (exact and
congestion-immune), while telemetry confirms service impact and is the path for
gray failures. Section~\ref{sec:detect} reports the empirical finding that drove this
split.

\paragraph{Resolution}
On a fault, a $k$-shortest-paths generator (Yen's algorithm over the live topology
graph, with the failed element removed) produces candidate survivors, and a
\emph{resolver} scores each with a scale-free multi-objective utility
\begin{equation}
\label{eq:obj}
U(p) = w_1(-\ell_p) + w_2\, c_p + w_3(-L_p) + w_4(-h_p),
\end{equation}
where $\ell_p$ is latency, $c_p$ capacity (bottleneck bandwidth), $L_p$ loss, and
$h_p$ hop count, each min--max normalized across candidates. The best path is
compiled to per-switch forward/reverse \texttt{OFPFlowMod} rules and installed
fail-secure. Recovery weights are fixed a~priori
($w=[0.5,1.0,1.5,0.1]$). % [VERIFIED docs/s2_5_results.txt]

\paragraph{Learning}
Each recovery's realized outcome (throughput, post-recovery loss, recovery time) is
scored by a reward $r=\beta\,\text{throughput}-\alpha\,\text{loss}-t_{\text{rec}}$
and stored. A \emph{learner} periodically re-fits the resolver weights from a recent
window of experience via the per-dimension covariance between each candidate's
normalized goodness and its reward, accepting an update only if a held-out
correlation does not drop (a forgetting guard). This lets SHIELD track
non-stationarity while keeping the resolver's objective interpretable.

\section{Evaluation}
We evaluate the three pillars on a seven-switch mesh with three node-disjoint
$s_1\!\leftrightarrow\!s_6$ paths (A:~5\,Mb/2\,ms, B:~10\,Mb/5\,ms, C:~20\,Mb/10\,ms)
and, for detection generalization, a second topology (\textsc{mesh2}). Faults are
injected by a scheduler (link/node down, congestion, flap); traffic is background
TCP+UDP with randomized bursts.

\subsection{Pillar 1 --- Detection and localization}
\label{sec:detect}
We first tested a pure-telemetry per-port classifier and found it \emph{cannot}
localize: a failed link and its innocent downstream neighbours look identical from
throughput, and under a shared data plane a link-down triggers network-wide
reconvergence rather than a localized signal---measured localization precision was
only $\approx0.10$. % [VERIFIED docs/p1_detection_results.txt]
This finding motivates the hybrid design. Using OpenFlow port-status for
localization and reserving telemetry for service-impact, SHIELD localizes
\emph{every} injected hard fault across both topologies with zero congestion false
alarms (Table~\ref{tab:detect}); because localization is deterministic port-status
rather than a learned model, it generalizes to the unseen \textsc{mesh2} by
construction.

\begin{table}[t]
\caption{Hybrid fault localization (every inter-switch link + a node fault + a
congestion window). % [VERIFIED docs/p1_detection_results.txt]}
\label{tab:detect}
\centering
\begin{tabular}{lcccc}
\toprule
Topology & Hard faults & Localized & Accuracy & Congestion FA \\
\midrule
\textsc{mesh}  & 10 & 10 & $1.000$ & 0 \\
\textsc{mesh2} &  9 &  9 & $1.000$ & 0 \\
\midrule
\textbf{Overall} & \textbf{19} & \textbf{19} & $\mathbf{1.000}$ & \textbf{0} \\
\bottomrule
\end{tabular}\\[2pt]
\footnotesize Mean localization latency $0.21$~s (poll-bounded; the production app is
event-driven).
\end{table}

\subsection{Pillar 2 --- Resolution against honest baselines}
\label{sec:resolve}
After the primary path (A) fails we measure the live QoS (iperf~TCP + ping) of each
surviving path under each scenario's injected conditions, then map every strategy to
its chosen path. SHIELD (which sees loss with observation noise $\sigma=0.05$) is
compared against \emph{fixed-failover} and \emph{shortest-path} (both loss-blind),
\emph{widest-path} (a strong load-aware baseline that routes on observed bottleneck
bandwidth), \emph{random} (control), and an \emph{oracle} that uses ground-truth
loss. Composite QoS $=\text{tput}-0.05\,\text{lat}-0.2\,\text{loss}$
(Table~\ref{tab:resolve}).

\begin{table}[t]
\caption{Composite QoS by strategy, 5 scenarios $\times$ 5 seeds (mean).
% [VERIFIED docs/s2_5_results.txt]}
\label{tab:resolve}
\centering
\footnotesize
\begin{tabular}{lrrrrrr}
\toprule
Scenario & fixed & short & widest & rand & \textbf{SHIELD} & oracle \\
\midrule
congest\_B      & $-0.2$ & $-0.2$ & $12.0$ & $2.2$  & $\mathbf{12.0}$ & $12.0$ \\
loss\_B         & $-4.0$ & $-4.0$ & $12.4$ & $-0.1$ & $\mathbf{12.4}$ & $12.4$ \\
congest\_loss\_B& $-6.9$ & $-6.9$ & $12.2$ & $-4.5$ & $\mathbf{12.2}$ & $12.2$ \\
congest\_C\_ctrl& $7.7$  & $7.7$  & $7.7$  & $5.3$  & $\mathbf{7.7}$  & $7.7$  \\
loss\_C\_fat    & $7.1$  & $7.1$  & $-11.7$& $4.0$  & $\mathbf{7.1}$  & $7.1$  \\
\bottomrule
\end{tabular}
\end{table}

SHIELD beats the loss-blind baselines in $3/5$ scenarios, is never materially worse
than widest-path, and \emph{strictly} beats it in \texttt{loss\_C\_fat}
($7.1$ vs $-11.7$), where 15\% loss lands on the high-capacity path so
bandwidth-only routing chooses the worst survivor. SHIELD matches the ground-truth
oracle in all five scenarios despite seeing only noisy loss, because the observed
loss gaps ($0$ vs $15$--$20\%$) exceed the noise. % [VERIFIED docs/s2_5_results.txt]
The honest reading: on pure-bandwidth faults load-awareness alone suffices (SHIELD
ties widest-path); the multi-objective, loss-aware policy earns its keep precisely
when loss decorrelates from bandwidth.

\subsection{Pillar 3 --- Learning on a non-stationary world}
\label{sec:learn}
We drive a nine-``day'' world whose degradation drifts: early days the loss lands on
Path~B (so C is best), late days it lands on Path~C (so B is best). Crucially the
degradation is \emph{loss only}---bandwidth stays nominal---and loss is the signal
established above to be unobservable at decision time, so the right path
\emph{cannot} be read from telemetry and must be learned from realized rewards.
QoS anchors are measured live on the mesh; the many-day, eight-seed replication
needed for a significance test is trace-driven from those measurements. We compare
three agents sharing the day-0 resolver: \textbf{SHIELD-online} (re-fits each day on
a two-day window, $\epsilon$-greedy), \textbf{frozen} (never updates), and
\textbf{adaptive} (a sliding-window-greedy non-stationary bandit---a principled
non-learning competitor that also tracks drift).

\begin{table}[t]
\caption{Mean recovery reward per day over 8 seeds
($r=\text{tput}-0.5\,\text{loss}-t_{\text{rec}}$).
% [VERIFIED docs/s3_4_results.txt]}
\label{tab:learn}
\centering
\footnotesize
\begin{tabular}{crrrc}
\toprule
Day & online & frozen & adaptive & $p(\text{degrade B})$ \\
\midrule
0 & $10.2$ & $13.0$  & $8.4$  & $0.90$ \\
2 & $8.7$  & $11.3$  & $9.1$  & $0.70$ \\
4 & $-1.6$ & $-1.7$  & $-1.7$ & $0.50$ \\
6 & $1.3$  & $0.3$   & $-0.5$ & $0.30$ \\
7 & $3.1$  & $-9.1$  & $1.0$  & $0.20$ \\
8 & $4.8$  & $-11.0$ & $6.7$  & $0.10$ \\
\bottomrule
\end{tabular}
\end{table}

As the world flips, the frozen model collapses from $13.0$ to $-11.0$ (the loss
migrates onto its fixed choice), while SHIELD-online and the adaptive baseline both
re-learn Path~B (Table~\ref{tab:learn}). Using an exact paired sign-flip permutation
test over eight seeds: in the late regime SHIELD-online exceeds frozen with
$p=0.0039$; the frozen model degrades early$\rightarrow$late with $p=0.0039$; the
online model shows no net catastrophic forgetting (held-out validation ends at
$0.46$ vs $0.32$ at day~0); and SHIELD-online \emph{matches} the adaptive baseline
($\text{online}=4.0$ vs $\text{adaptive}=3.8$; $p=0.547$ that adaptive beats online).
% [VERIFIED docs/s3_4_results.txt]
The claim we defend is therefore the honest one: SHIELD adapts day-by-day as well as
a purpose-built adaptive bandit, but does so by updating an \emph{interpretable}
multi-objective resolver rather than an opaque arm table.

\section{Discussion and Threats to Validity}
\label{sec:threats}
\textbf{Where a simpler method suffices.} We are explicit that the contribution is
the integrated framework, not a learning algorithm. With accurate telemetry the
heuristic resolver already matches the oracle (Table~\ref{tab:resolve}), and the
learner's value appears only under partial observability and drift
(Table~\ref{tab:learn}). Reinforcement learning is retained as a component and
ablation, not the headline.

\textbf{Emulation and scale.} Results are on Mininet/OVS at a scale sufficient for a
controlled study; the resolution study uses two survivors after a single primary
failure, and the learning study is a two-arm non-stationary bandit with a binary
decisive feature. Extending the resolver and learner to $\geq3$-path topologies and
richer feature interactions is the primary generalization threat and is future work.

\textbf{Trace-driven replication.} The learning curve draws outcomes from
live-measured QoS anchors rather than a fully-live nine-day loop; the anchors are
real but the day-scale replication is simulated to obtain a $p$-value. A fully-live
multi-day run is future work.

\textbf{Statistical scope.} We use exact paired sign-flip permutation tests
(stdlib, no external solver). Eight seeds give a comfortable margin under $0.05$
($p=0.0039$) with no seed on the boundary; a pre-registered sensitivity grid over the
window/exploration hyper-parameters, and a validation hold-out drawn from a
\emph{prior} regime (to probe forgetting of old regimes), would further harden the
learning claim.

\section{Conclusion and Future Work}
SHIELD integrates fault detection, path resolution, and online learning into one
self-healing SDN control loop and evaluates each stage against an honest baseline: a
ground-truth oracle and a load-aware widest-path router for resolution, and a frozen
ablation plus an adaptive bandit for learning. It localizes every injected hard
fault, resolves faults at least as well as a strong load-aware baseline and better
when loss decorrelates from bandwidth, and adapts to a drifting world as well as a
purpose-built adaptive competitor---without catastrophic forgetting. Future work
targets larger multi-path topologies, gray-failure detection through telemetry,
a fully-live multi-day learning loop, and a pre-registered hyper-parameter
sensitivity study. All controllers, topologies, the fault-injection harness, and
the evaluation and reproduction scripts are released.

\bibliographystyle{IEEEtran}
\bibliography{references}

\end{document}
